%!TEX root = EM_singular_models.tex

% PDF margin etc settings
%\setlength{\textwidth}{\paperwidth}
%\addtolength{\textwidth}{-6cm}
%\setlength{\textheight}{\paperheight}
%\addtolength{\textheight}{-4cm}
%\addtolength{\textheight}{-1.1\headheight}
%\addtolength{\textheight}{-\headsep}
%\addtolength{\textheight}{-\footskip}
%\setlength{\oddsidemargin}{0.5cm}
%\setlength{\evensidemargin}{0.5cm}

%%%%%%%%%%%%%%%%%%%%%%%%%%%%%%%%

% MACROS HERE

%%%%%%%%%%%%%%%%%%%%%%%%%%%%%%%%

\newcommand{\fixedpointtheta}{\widehat\theta_\obs}

\newcommand{\cover}{N}
\newcommand{\kinit}{\nu}
\newcommand{\cdelta}{c_\threshold}
\newcommand{\dndelta}{\omega}
\newcommand{\kndelta}{\Upsilon}
\newcommand{\kndeltanew}{\Psi}
\newcommand{\radii}{\mathcal{R}}
\newcommand{\event}{\mathcal{E}}
\newcommand{\kevent}{\mathcal{D}}


% Observations, dimension etc.
\newcommand{\dims}{\ensuremath{d}}
\newcommand{\samples}{\ensuremath{n}}
\newcommand{\obs}{\samples}
\newcommand{\real}{\ensuremath{\mathbb{R}}}
\newcommand{\complex}{\ensuremath{\mathbb{C}}}
\newcommand{\realdim}{\ensuremath{\real^\dims}}

\newcommand{\smallthreshold}{\epsilon}

% some mathcal notations
\newcommand{\order}[1]{\ensuremath{\mathcal{O}\parenth{#1}}}

% Basic statistics notation
\newcommand{\xstar}{\ensuremath{x^*}}
\newcommand{\xbar}{\ensuremath{\bar{x}}}
% True parameter
\newcommand{\thetastar}{\ensuremath{{\theta^*}}}
% Estimate one
\newcommand{\thetahat}{\ensuremath{\widehat{\theta}}}
% Estimate two
\newcommand{\thetatil}{\ensuremath{\widetilde{\theta}}}
\newcommand{\thetabar}{\ensuremath{\bar{\theta}}}
\newcommand{\yvec}{\ensuremath{y}}
\newcommand{\Xmat}{\ensuremath{X}}

% Distributions
\newcommand{\distribution}{P}
\newcommand{\density}{f}
\newcommand{\gaussiandensity}{\phi}
\newcommand{\NORMAL}{\ensuremath{\mathcal{N}}}
\newcommand{\BER}{\ensuremath{\mbox{Ber}}}

% Spaces
\newcommand{\Xspace}{\ensuremath{\mathcal{X}}}
\newcommand{\Yspace}{\ensuremath{\mathcal{Y}}}
\newcommand{\Zspace}{\ensuremath{\mathcal{Z}}}


% Brackets Size
\newcommand{\brackets}[1]{\left[ #1 \right]}
\newcommand{\parenth}[1]{\left( #1 \right)}
\newcommand{\bigparenth}[1]{\big( #1 \big)}
\newcommand{\biggparenth}[1]{\bigg( #1 \bigg)}
\newcommand{\braces}[1]{\left\{ #1 \right \}}
\newcommand{\abss}[1]{\left| #1 \right |}
\newcommand{\angles}[1]{\left\langle #1 \right \rangle}
\newcommand{\tp}{^\top}

% Generic vectors and scalars
\newcommand{\myvec}{u} % for generic vector
\newcommand{\myvectwo}{v} % for generic vector
\newcommand{\mymat}{B}
\newcommand{\set}{{A}}
\newcommand{\scalar}{\alpha}
\newcommand{\tvscalar}{\epsilon}
\newcommand{\tailboundscalar}{t}
\newcommand{\term}{U}
\newcommand{\myVarone}{\alpha_X}
\newcommand{\myVartwo}{\alpha_Y}
\newcommand{\myCovmat}{\Sigma}
\newcommand{\xcord}[1]{x^{(#1)}}
\newcommand{\xcordpop}[1]{X^{(#1)}}
\newcommand{\ycordpop}[1]{Y^{(#1)}}
\newcommand{\myVarthree}{\ensuremath{U}}
\newcommand{\myVarfour}{\ensuremath{V}}
\newcommand{\myfuntwo}{\ensuremath{g}}
\newcommand{\numsteps}{t}
\newcommand{\totalnumsteps}{T}
\newcommand{\seqalpha}[1]{\alpha_{#1}}
\newcommand{\seqnu}[1]{\nu^{(#1)}}
\newcommand{\imax}{\ell_\smallthreshold}
\newcommand{\basis}{e}
\newcommand{\radem}{\varepsilon}
\newcommand{\Tone}{T\ensuremath{_{0}}}
\newcommand{\Ttwo}{T\ensuremath{_{2}}}
\newcommand{\contractPar}[1]{\ensuremath{\gamma\parenth{#1}}}

\newcommand{\simiid}{\overset{\mathrm{i.i.d.}}{\sim}}
\newcommand{\eqd}{\overset{d}{=}}


% EM updates
\newcommand{\hidden}{Z}
\newcommand{\qfun}[2]{Q(#1\vert#2)}
\newcommand{\variable}{\theta}
\newcommand{\weights}{w}
\newcommand{\myfun}{q}
\newcommand{\gradf}{\nabla \myfun}
\newcommand{\hessf}{\nabla^2 \myfun}
\newcommand{\iter}{t}
\newcommand{\mupdate}{M}
\newcommand{\mupdategeneral}{\ensuremath{M^{GN}}}




% EM contractions
\newcommand{\contractgene}{\ensuremath{\gamma^{GN}}}
\newcommand{\updateMat}{\ensuremath{\Gamma_{\theta}}}
\newcommand{\mymatTheta}{\ensuremath{B_{\theta}}}
\newcommand{\sech}{\text{sech}}

% Nhat's macros 
\newcommand{\Rspace}{\ensuremath{\mathbb{R}}}
\newcommand{\Ecal}{\ensuremath{\mathcal{E}}}
\newcommand{\Ncal}{\ensuremath{\mathcal{N}}}
\newcommand{\gradient}{\ensuremath{\nabla}}
\newcommand{\smallweight}{\ensuremath{\epsilon^{sw}}}
\newcommand{\smallnorm}{\ensuremath{\rho^{sw}}}
\newcommand{\ball}{\ensuremath{\mathbb{B}}}
\newcommand{\sphere}{\ensuremath{\mathbb{S}}}

\newcommand{\diam}{\ensuremath{\text{diam}}}
\newcommand{\upper}{\ensuremath{\bar{C}}}
\newcommand{\sample}{\ensuremath{\xi}}
\newcommand{\weight}{\ensuremath{\pi}}
\newcommand{\contracasym}{\ensuremath{\rho}}
\newcommand{\Fishermatrix}{\ensuremath{\mathcal{I}}}
\newcommand{\barcontracasym}{\ensuremath{\bar{\rho}}}


\newcommand{\TrueDensity}{\ensuremath{f}}
\newcommand{\FitDensity}{\ensuremath{f}_{\theta}}
\newcommand{\FitDensityfree}{\ensuremath{f}_{\pi, \theta}}
\newcommand{\FitDensitygen}{\ensuremath{f}_{\pi, \theta_{1}, \theta_{2}}}
\newcommand{\FitDensitymore}{\ensuremath{f}_{G}}
\newcommand{\normden}{\ensuremath{\phi}}
\newcommand{\mean}{\ensuremath{\theta}}
\newcommand{\sd}{\ensuremath{\sigma}}
\newcommand{\parFOS}{\ensuremath{\gamma}}
\newcommand{\funQ}{\ensuremath{Q}}
\newcommand{\updateM}{\ensuremath{M}}
\newcommand{\weightFun}{\ensuremath{w}}
\newcommand{\mydefn}{\ensuremath{:=}}
\newcommand{\poincare}{\ensuremath{\tilde{C}}}
\newcommand{\paraspace}{\ensuremath{\Theta}}
\newcommand{\loglihood}{\ensuremath{F}}
\newcommand{\prior}{\ensuremath{\pi}}
\newcommand{\posterior}{\ensuremath{\Pi}}
\newcommand{\boundcons}{\ensuremath{B}}
\newcommand{\boundconstot}{\ensuremath{H}}
\newcommand{\noise}{\ensuremath{\varepsilon}}
\newcommand{\loglihoodgaus}{\ensuremath{F^{G}}}
\newcommand{\loglihoodind}{\ensuremath{F^{I}}}
\newcommand{\loglihoodlogit}{\ensuremath{F^{R}}}
\newcommand{\noiseone}{\ensuremath{\varepsilon_{1}}}
\newcommand{\noisetwo}{\ensuremath{\varepsilon_{2}}}

%%%SDE_Bayesian_inf
\newcommand{\weakcon}{\ensuremath{\psi}}
\newcommand{\perturb}{\ensuremath{\zeta}}
\newcommand{\inver}{\ensuremath{\xi}}

% Universal constants
\newcommand{\unicon}{c}
\newcommand{\uniconnew}{c'}
\newcommand{\uniconone}{c_1}
\newcommand{\unicontwo}{c_2}

% some parameters that you may define
\newcommand{\scparam}{m} % strong  convexity parameter
\newcommand{\smoothness}{L} % smoothness parameter
\newcommand{\step}{h}


%%%%%%%%% Distributions and Random variables %%%%%%%%%%%
\newcommand{\rvg}{\xi} % to denote the random variable g
\newcommand{\threshold}{\delta}


%%%%%%%%% Basic Terms like defn, etal, tmix, polylog %%%%%%%%%%%
\newcommand{\defn}{:=}
\newcommand{\rdefn}{=:}
\newcommand{\etal}{{et al.}}
\newcommand{\polylog}{\text{poly-log}}

% Some vector/matrix norms
\newcommand{\matnorm}[3]{|\!|\!| #1 | \! | \!|_{{#2}, {#3}}}
\newcommand{\matsnorm}[2]{|\!|\!| #1 | \! | \!|_{{#2}}}
\newcommand{\vecnorm}[2]{\left\| #1\right\|_{#2}}
\newcommand{\enorm}[1]{\vecnorm{#1}{2}} % euclidean norm
\newcommand{\nucnorm}[1]{\ensuremath{\matsnorm{#1}{\footnotesize{\mbox{nuc}}}}}
\newcommand{\fronorm}[1]{\ensuremath{\matsnorm{#1}{\footnotesize{\mbox{fro}}}}}
\newcommand{\opnorm}[1]{\ensuremath{\matsnorm{#1}{\tiny{\mbox{op}}}}}
%%% EM paper
\newcommand{\constrainnorm}[1]{\vecnorm{#1}{[k]}} % euclidean norm

\newcommand{\bigmatsnorm}[2]{{\Big|\!\Big|\!\Big| #1
  \Big|\!\Big|\!\Big|}_{#2}}

% Inner product
\newcommand{\inprod}[2]{\ensuremath{\langle #1 , \, #2 \rangle}}

% Kullback-Leibler
\newcommand{\kull}[2]{\ensuremath{D_{\text{KL}}(#1\; \| \; #2)}}

% Probability
\newcommand{\Exs}{\ensuremath{{\mathbb{E}}}}
\newcommand{\Prob}{\ensuremath{{\mathbb{P}}}}

%Eigenvector / eigenvalue related notation
\newcommand{\eig}[1]{\ensuremath{\lambda_{#1}}}
\newcommand{\eigmax}{\ensuremath{\eig{\max}}}
\newcommand{\eigmin}{\ensuremath{\eig{\min}}}

% \DeclareMathOperator{\det}{det}
\DeclareMathOperator{\modd}{mod}
\DeclareMathOperator{\diag}{diag}
\DeclareMathOperator{\var}{var}
\DeclareMathOperator{\cov}{cov}
\DeclareMathOperator{\trace}{trace}
\DeclareMathOperator{\abs}{abs}
\DeclareMathOperator{\floor}{floor}
\DeclareMathOperator{\vol}{vol}
\DeclareMathOperator{\child}{child}
\DeclareMathOperator{\parent}{parent}
\DeclareMathOperator{\sign}{sign}
\DeclareMathOperator{\rank}{rank}
\DeclareMathOperator{\toeplitz}{toeplitz}




\newtheoremstyle{named}{}{}{\itshape}{}{\bfseries}{.}{.5em}{\thmnote{#3's }#1}
\theoremstyle{named}
\newtheorem*{namedtheorem}{Lemma}

%%%%%%%
\theoremstyle{plain}

% {Theorem, Proposition, Lemma, Corollary} numbered sequentially
% throughout the paper
\newtheorem{theorem}{Theorem}
\newtheorem{proposition}{Proposition}
\newtheorem{lemma}{Lemma}
\newtheorem{claim}{Claim}
\newtheorem{corollary}{Corollary}
\newtheorem{definition}{Definition}
\newtheorem{conjecture}{Conjecture}
\newtheorem{remark}{Remark}
%%%%%%%%%%%%%%%%%%%%%%%%%%%%%%%%%%%%%%%%%%%%%%%%%%%%%%%%%%%%%%%%%%%%%%%
% WIDEBAR COMMAND
\newlength{\widebarargwidth}
\newlength{\widebarargheight}
\newlength{\widebarargdepth}
\DeclareRobustCommand{\widebar}[1]{%
  \settowidth{\widebarargwidth}{\ensuremath{#1}}%
  \settoheight{\widebarargheight}{\ensuremath{#1}}%
  \settodepth{\widebarargdepth}{\ensuremath{#1}}%
  \addtolength{\widebarargwidth}{-0.3\widebarargheight}%
  \addtolength{\widebarargwidth}{-0.3\widebarargdepth}%
  \makebox[0pt][l]{\hspace{0.3\widebarargheight}%
    \hspace{0.3\widebarargdepth}%
    \addtolength{\widebarargheight}{0.3ex}%
    \rule[\widebarargheight]{0.95\widebarargwidth}{0.1ex}}%
  {#1}}



%%% New version of \caption puts things in smaller type, single-spaced
%%% and indents them to set them off more from the text.
\makeatletter
\long\def\@makecaption#1#2{
        \vskip 0.8ex
        \setbox\@tempboxa\hbox{\small {\bf #1:} #2}
        \parindent 1.5em  %% How can we use the global value of this???
        \dimen0=\hsize
        \advance\dimen0 by -3em
        \ifdim \wd\@tempboxa >\dimen0
                \hbox to \hsize{
                        \parindent 0em
                        \hfil
                        \parbox{\dimen0}{\def\baselinestretch{0.96}\small
                                {\bf #1.} #2
                                %%\unhbox\@tempboxa
                                }
                        \hfil}
        \else \hbox to \hsize{\hfil \box\@tempboxa \hfil}
        \fi
        }
\makeatother



%% COMMENTING commands

\long\def\comment#1{}
\definecolor{battleshipgrey}{rgb}{0.52, 0.52, 0.51}
\definecolor{darkgray}{rgb}{0.66, 0.66, 0.66}
\definecolor{darkgreen}{rgb}{0.0, 0.2, 0.13}
\definecolor{darkspringgreen}{rgb}{0.09, 0.45, 0.27}
\definecolor{dukeblue}{rgb}{0.0, 0.0, 0.61}
\definecolor{olivedrab7}{rgb}{0.24, 0.2, 0.12}
\definecolor{darkblue}{rgb}{0.0, 0.0, 0.55}
\definecolor{darkscarlet}{rgb}{0.34, 0.01, 0.1}
\definecolor{candyapplered}{rgb}{1.0, 0.03, 0.0}
\definecolor{ao(english)}{rgb}{0.0, 0.5, 0.0}
\definecolor{applegreen}{rgb}{0.55, 0.71, 0.0}

\newcommand{\red}[1]{\textcolor{red}{#1}}
\newcommand{\mjwcomment}[1]{{\bf{{\red{{MJW --- #1}}}}}}
\newcommand{\mwlcomment}[1]{{\bf{{\color{orange}{{Wenlong --- #1}}}}}}

\newcommand{\blue}[1]{\textcolor{blue}{#1}}
\newcommand{\green}[1]{\textcolor{green}{#1}}
\newcommand{\highlight}[1]{\textcolor{applegreen}{#1}}

% comment lines
\newcommand{\genecomment}[1]{{\bf{{\blue{{TEAM --- #1}}}}}}
\newcommand{\rrdcomment}[1]{{\bf{{\blue{{RRD --- #1}}}}}}
\newcommand{\nhcomment}[1]{{\bf{{\blue{{NH --- #1}}}}}}
\newcommand{\kkcomment}[1]{{\bf{{\darkgreen{{KK --- #1}}}}}}


\newcommand{\widgraph}[2]{\includegraphics[keepaspectratio,width=#1]{#2}}
